\section{Background of the Study}

Video compression for surveillance sits at the intersection of three established bodies of knowledge: general video coding, content-aware (ROI-based) compression, and edge-based video analytics. General video coding is defined by the H.264/AVC standard \cite{Wiegand2003H264} and its successor HEVC \cite{Sullivan2012HEVC}, which together underpin the overwhelming majority of deployed CCTV encoders. Content-aware compression modifies or extends these codecs so that semantically important regions receive preferential quality \cite{Huang2009MultiPriority,Xue2016FastROI}. Edge-based analytics moves scene understanding close to the camera so that transmission and storage decisions can be informed by content \cite{Chen2019DIVS,Nguyen2020Scalable}. This chapter reviews 45 references organized thematically, rather than as a sequential list of papers, and closes by identifying the precise research gap addressed by this thesis. All claims about the content of cited works are drawn from the bibliographic evidence available for each reference (title, venue, and supplied metadata); where only metadata was available, no methodological detail beyond that metadata is asserted.

\section{Related Work / Existing Systems}

\subsection{Video Coding Standards and Early Surveillance Compression}

The baseline against which any surveillance compression method must be judged is uniform encoding with a standard codec. Wiegand \textit{et al.} \cite{Wiegand2003H264} document the H.264/AVC standard, and Sullivan \textit{et al.} \cite{Sullivan2012HEVC} document HEVC, which approximately doubles the coding efficiency of its predecessor at the cost of increased computational complexity. Both standards support quantization-parameter (QP) variation across regions, which is the mechanism exploited by most ROI-based extensions. Early work on adaptive video compression specifically for surveillance applications \cite{ISM2011Adaptive} established the basic premise that surveillance sequences---characterized by static backgrounds and sporadic foreground activity---are poorly served by content-blind encoding, a premise confirmed at scale by studies of smart monitoring cameras and big surveillance video data \cite{Shao2018Smart}.

\subsection{Classical and Heuristic ROI-Based Coding}

A first generation of ROI methods used hand-crafted priority maps and foveation models. Huang and Lin \cite{Huang2009MultiPriority} proposed multiple-priority ROI H.264 compression with constrained variable bitrate control, demonstrating that distinct priority levels can be enforced within a standard-compliant bitstream. Galan-Hernandez \textit{et al.} \cite{Galan2011Foveated} applied foveated compression with hierarchical trees for real-time transmission, trading peripheral fidelity for foreground quality in a biologically inspired manner. A codec-independent ROI framework using joint pre- and post-processing \cite{CodecIndependent2016} showed that ROI effects can alternatively be achieved outside the codec itself, by degrading the background before encoding and restoring structure after decoding---an approach conceptually close to the background downscaling path implemented in this thesis. Xue and Wei \cite{Xue2016FastROI} addressed the computational bottleneck directly with fast ROI-based HEVC coding for surveillance video, recognizing that ROI methods must run within the real-time budget of the encoder to be deployable.

\subsection{Object-Aware and Foreground/Background Compression}

A second strand makes the ROI object-aware: regions are defined by detected objects or by foreground segmentation rather than by fixed geometry. Xiao \textit{et al.} \cite{Xiao2017Sensitive} proposed a sensitive-object-oriented approach to surveillance compression for smart cities, prioritizing socially and legally sensitive content. Ma \textit{et al.} \cite{Ma2016Lossless} studied lossless ROI privacy protection in H.264/AVC streams, showing that ROI machinery can serve privacy as well as efficiency---a duality reflected in this thesis through its privacy modes. Meuel \textit{et al.} \cite{Meuel2018Aerial} evaluated ROI coding for aerial surveillance with AVC and HEVC, confirming that the ROI advantage depends on scene characteristics such as camera motion. Wu \textit{et al.} \cite{Wu2020FGBG} introduced foreground-background parallel compression with residual encoding, reporting bitrate reductions of approximately 49.75\% in bits per pixel against H.265 at 36 dB on their test data; this method is directly relevant because one of the baselines evaluated in this thesis (foreground/background split) is an instance of the same family, and its failure on moving-camera content (Chapter~7) illustrates the scene-dependence Meuel \textit{et al.} warned about. Hapsari \textit{et al.} \cite{Hapsari2022ROI} further reported ROI-based gains on surveillance systems, consistent with the broader pattern that static-camera sequences benefit most.

\subsection{Edge-Cloud Collaborative Video Analytics}

The architectural context of this thesis---detection at the edge, management and visualization in the cloud or on a server---is supported by a substantial literature. Chen \textit{et al.} \cite{Chen2019DIVS} presented a distributed deep-learning model for intelligent surveillance with edge computing. Elgamal \textit{et al.} \cite{Elgamal2020SIEVE} proposed semantically encoded video analytics across edge and cloud, encoding semantic information so that downstream analytics can operate at reduced bitrates. Nguyen and Choi \cite{Nguyen2020Scalable} showed that compressed-domain features can be exploited at the edge for scalable analytics, reducing the need to fully decode streams before analysis. Li \textit{et al.} \cite{Li2022EdgeComp} combined surveillance video analysis with compression explicitly on edge-computing platforms. Wang \textit{et al.} \cite{Wang2022E2ESurv} studied end-to-end compression tuned for surveillance and conference videos rather than for generic content. At larger scale, Nan \textit{et al.} \cite{Nan2023Collaborative} addressed cloud-edge collaborative continuous learning for large-scale analytics, and Zhang \textit{et al.} \cite{Zhang2022JCAB} studied joint configuration selection and bandwidth allocation for edge-based analytics---the resource-allocation problem of which this thesis's Lagrangian-inspired allocator is a small, heuristic instance.

\subsection{Recognition-Aware and Adaptive Allocation}

Several studies allocate bits to preserve machine-vision performance rather than human-viewed fidelity. Beye \textit{et al.} \cite{Beye2023Recognition} proposed recognition-aware bitrate allocation for AI-enabled remote surveillance, directly coupling encoding decisions to detector accuracy. Wang and Yang \cite{Wang2025Metadata} used compression-metadata-assisted ROI extraction and adaptive inference, showing that the codec and the detector can inform each other bidirectionally. Yang \textit{et al.} \cite{Yang2025RegionAware} proposed adaptive region-aware encoding for real-time cloud-edge collaborative detection. Qian \textit{et al.} \cite{Qian2025OsmoticGate} applied multi-agent deep reinforcement learning to edge-cloud collaborative streaming analytics. Perugachi-Diaz \textit{et al.} \cite{Perugachi2022NeuralROI} and Li \textit{et al.} \cite{Li2023SwinROI} investigated neural and transformer-based ROI compression, while a VANET-based study \cite{VANET2023ROI} demonstrated ROI-adaptive cross-layer video streaming under volatile wireless conditions---evidence that ROI control generalizes beyond wired CCTV. The video-coding-for-machines paradigm \cite{Duan2020VCM} provides the theoretical umbrella for this entire direction: when the consumer of video is an algorithm, rate--distortion optimization should be redefined around task performance.

\subsection{Learned and Recent Adaptive Surveillance Coding}

Recent learned codecs have pushed surveillance compression further under controlled conditions. Yang \textit{et al.} \cite{Yang2024AVCSR} proposed adaptive compressive sensing with ROI detection in the compressed domain. Wu \textit{et al.} \cite{Wu2024ROIDVC} introduced ROI-DVC, a deep video-coding framework built around regions of interest, reporting approximately 32\% bitrate reduction against H.265 at matched quality on their data. Wu \textit{et al.} \cite{Wu2025How2Compress} proposed adaptive granular compression for scalable edge analytics, reporting up to approximately 50.4\% reduction (up to 3.01$\times$ in some configurations) with accuracy held. Zhao \textit{et al.} proposed background-hyperprior \cite{Zhao2025ASVC} and temporal-adaptive learned \cite{Zhao2025TALC} surveillance codecs. Wang \textit{et al.} \cite{Wang2025CAVC} studied content-aware compression for roadside cameras; Pu \textit{et al.} \cite{Pu2025AdaptiveBitrate} proposed adaptive ROI bitrate encoding driven by visual recognition; Chuang and Hsiao \cite{Chuang2025AdaptiveROI} combined adaptive ROI encoding with caching for surveillance streaming; PAVC \cite{PAVC2025} introduced precision-aware compression for bandwidth reduction; Baid \textit{et al.} \cite{Baid2026BLUE} addressed semantics-preserving compression for vision-language analytics; and Kadukar and Dani \cite{Kadukar2026Semantic} reported approximately 74.63\% data reduction with a semantic-aware offloading framework. VaBUS \cite{Wang2023VaBUS} demonstrated edge-cloud real-time analytics via background understanding and subtraction, reporting bandwidth reductions in the range of approximately 25.0--76.9\%. These figures are cited here strictly as the results reported by their respective authors on their own datasets and protocols; they are not re-measured in this thesis and are not claimed to transfer to the present deployment.

\subsection{Forensic Integrity and Chain of Evidence}

Surveillance video is frequently evidentiary, so integrity mechanisms matter alongside compression. Kanwal \textit{et al.} \cite{Kanwal2020Chain} studied the preservation of chain-of-evidence in surveillance videos for authentication and trust-enabled sharing. Their work frames the requirement that motivated the forensic chain in this thesis: every stored event frame should be accompanied by verifiable integrity metadata. The present implementation realizes a narrow, exact-byte form of that vision using SHA-256 hash chains over JPEG event frames and JSON sidecars (Chapter~6), and the evaluation (Chapter~7) tests precisely what such a chain can and cannot detect.

\subsection{Anomaly and Event-Aware Surveillance}

Because this thesis probes its risk score against anomaly-labelled data (UCF-Crime), the anomaly-detection literature is relevant as context---and as a warning. Duong \textit{et al.} \cite{Duong2023Survey} survey deep-learning-based anomaly detection in video surveillance; a systematic ACM taxonomy \cite{ACMTaxonomy2024} compares deep models for generalized anomaly event detection; Baradaran and Bergevin \cite{Baradaran2021Anomaly} critically examine semi-supervised anomaly-detection methods; and autoencoder-based abnormal-event detection systems \cite{Autoencoder2024Event} exemplify the reconstruction-error approach. The consistent message of these works is that anomaly labels denote rare, context-defined events, whereas a heuristic object-presence score denotes something far narrower. That distinction is central to the interpretation of Phase~1 results in Chapter~7.

\section{Research Gap}

The studies reviewed above primarily focus on (i) ROI constructions evaluated on static-camera or otherwise favourable sequences, (ii) learned codecs requiring GPU-class resources, (iii) recognition-aware allocation evaluated for analytic accuracy rather than joint compression--integrity--visualization deployment, and (iv) forensic or anomaly capabilities studied separately from compression. Fewer studies address a complete, CPU-deployable edge-cloud CCTV prototype that combines heuristic ROI control, state-driven adaptive encoding, tamper-evident event buffering, and dashboard visualization, and evaluates all of them under one reproducible, honestly reported protocol that includes negative results. In particular, controlled same-codec comparisons with shared detection caches, explicit BD-rate withholding criteria, group-aware calibration with degeneracy analysis, and pre-registered forensic expectations are rarely reported together. This thesis occupies that gap: not by claiming a new coding-theoretic optimum, but by building the integrated prototype and documenting---quantitatively and without embellishment---where it works, where it fails, and why.

\section{Research Questions}

\begin{itemize}
\item RQ1: Can dynamic ROI-based control driven by YOLOv8n be implemented in a practical edge-based CCTV compression pipeline with state-dependent encoding and event buffering?
\item RQ2: How does the implemented heuristic risk/ROI mechanism behave across public video datasets (CDnet-2014, UCF-Crime), and what is the precise relationship between the risk score and anomaly labels?
\item RQ3: How does the proposed ROI allocation compare with uniform encoding and alternative ROI baselines at matched bitrate under a locked same-codec protocol?
\item RQ4: Can the event-buffering mechanism provide tamper-evident integrity verification under controlled tampering, recompression, and clean-control conditions, and what are the boundaries of that guarantee?
\end{itemize}

RQ1 is addressed by Chapters~5--6 (construction evidence); RQ2 by Phases~1--2 (Chapters~4 and 7); RQ3 by Phase~4; and RQ4 by Phase~3. The corresponding evaluation protocols are defined in Chapter~4, and the results with their limitations are reported in Chapter~7.
